\section*{Организация процесса тестирования}

Для организации процесса тестирования и контроля выполнения задач использовалась облачная система управления проектами Kaiten. Было создано рабочее пространство и доска проекта. Задачи, подлежащие тестированию, зафиксированы в виде карточек. 

Для каждой карточки заданы следующие атрибуты: статус выполнения, дата начала, дата окончания, пользовательская метка. График выполнения задач в виде доски представлен на \cref{fig:kaiten_table}.

\begin{figure}[H]
    \centering
    \includegraphics[width=\textwidth]{figs/kaiten_table.png}
    \caption{График выполнения задач в виде доски}
    \label{fig:kaiten_table}
\end{figure}

Диаграмма Ганта, отображающая длительность выполнения каждой задачи и временные рамки, представлена на \cref{fig:kaiten_timeline}. График выполнения задач в табличном виде показан на \cref{fig:kaiten_report}.

\begin{figure}[H]
    \centering
    \includegraphics[width=\textwidth]{figs/kaiten_timeline.png}
    \caption{TIMELINE-диаграмма выполнения проекта}
    \label{fig:kaiten_timeline}
\end{figure}

\begin{figure}[H]
    \centering
    \includegraphics[width=\textwidth]{figs/kaiten_report.png}
    \caption{График выполнения задач в табличном виде}
    \label{fig:kaiten_report}
\end{figure}

\section{Разработка тестовых мероприятий}

В качестве объекта тестирования выбрано веб-приложение Ассоциации \textit{Цифровая Энергетика} (URL: \url{https://www.digital-energy.ru/}). Разработан план тестовых мероприятий, включающий проверки графического интерфейса, совместимости, адаптивности и сетевого взаимодействия. План зафиксирован в \cref{tab:test_plan}.

\begin{longtblr}[
    caption = {План тестовых мероприятий},
    label = {tab:test_plan},
]{
    colspec = { Q[l,m] X[l,m] X[l,m] Q[l,m] },
    hlines, vlines,
    rowhead = 1,
}
    № & Описание проверки & Ожидаемый результат & Статус \\
    1 & Проверка загрузки главной страницы & Страница загружается без ошибок, код 200 & Успешно \\
    2 & Проверка отображения логотипа в шапке & Логотип отображается, является ссылкой на главную & Успешно \\
    3 & Работоспособность главного меню навигации & При клике происходит переход в соответствующие разделы & Успешно \\
    4 & Тестирование совместимости: Windows + Firefox & Корректное отображение элементов, отсутствие смещений & Успешно \\
    5 & Тестирование совместимости: Windows + Opera & Корректное отображение элементов, отсутствие смещений & Успешно \\
    6 & Тестирование совместимости: Windows + Vivaldi & Корректное отображение элементов, отсутствие смещений & Успешно \\
    7 & Тестирование адаптивности: разрешение $1920\times1080$ & Десктопная версия отображается корректно & Успешно \\
    8 & Тестирование адаптивности: разрешение $1366\times768$ & Версия для ноутбуков отображается без горизонтальной прокрутки & Успешно \\
    9 & Тестирование адаптивности: разрешение $360\times800$ & Мобильная версия активируется, появляется гамбургер-меню & Успешно \\
    10 & Проверка ссылок в подвале сайта (footer) & Все ссылки ведут на существующие страницы (нет 404) & Успешно \\
    11 & Валидация HTML-кода главной страницы (W3C) & Отсутствие критических ошибок верстки & Замечания \\
    12 & Проверка наличия favicon & Иконка вкладки отображается в браузере & Успешно \\
    13 & Проверка работы формы обратной связи & Данные отправляются, выводится сообщение об успехе & Успешно \\
    14 & Валидация полей формы (пустые поля) & Форма не отправляется, выводится предупреждение & Успешно \\
    15 & Проверка скорости загрузки элементов (Network) & Время загрузки DOM не превышает допустимых пределов & Успешно \\
    16 & Проверка структуры заголовков (H1-H6) & Иерархия соблюдена, на странице один H1 & Успешно \\
    17 & Проверка наличия атрибутов alt у изображений & Все смысловые изображения имеют текстовое описание & Замечания \\
    18 & Технический аудит используемых технологий & Определен стек технологий веб-приложения & Успешно \\
    19 & Анализ заголовков HTTP-ответов & Присутствуют заголовки безопасности & Успешно \\
    20 & Проверка файлов cookie & Файлы cookie имеют флаги Secure и HttpOnly & Успешно \\
\end{longtblr}

\section{Тестирование совместимости}

Тестирование совместимости проводилось методом попарного тестирования различных браузеров в среде ОС Windows 10. Цель — убедиться в кроссбраузерной идентичности отображения контента. Результаты проверки в браузерах Firefox 133, Opera 108 и Vivaldi 7.0 представлены на \cref{fig:compat_firefox,fig:compat_opera,fig:compat_vivaldi}.

\begin{figure}[H]
    \centering
    \includegraphics[width=\textwidth]{figs/firefox1.png}
    \caption{Отображение веб-приложения в браузере Firefox 133}
    \label{fig:compat_firefox}
\end{figure}

\begin{figure}[H]
    \centering
    \includegraphics[width=\textwidth]{figs/opera1.png}
    \caption{Отображение веб-приложения в браузере Opera 108}
    \label{fig:compat_opera}
\end{figure}

\begin{figure}[H]
    \centering
    \includegraphics[width=\textwidth]{figs/vivaldi1.png}
    \caption{Отображение веб-приложения в браузере Vivaldi 7.0}
    \label{fig:compat_vivaldi}
\end{figure}

Расхождений не наблюдается.

\section{Тестирование адаптивности}

Для проведения тестирования адаптивности использовались встроенные инструментальные средства разработчика DevTools. Проверка поведения графического интерфейса выполнена на трех целевых разрешениях экрана: $1920\times1080$, $1366\times768$ и $360\times800$. Результаты зафиксированы на \cref{fig:adapt_1920,fig:adapt_1366,fig:adapt_360}.

\begin{figure}[H]
    \centering
    \includegraphics[width=\textwidth]{figs/1920+1080.png}
    \caption{Тестирование адаптивности для разрешения $1920\times1080$}
    \label{fig:adapt_1920}
\end{figure}

\begin{figure}[H]
    \centering
    \includegraphics[width=\textwidth]{figs/1366+768.png}
    \caption{Тестирование адаптивности для разрешения $1366\times768$}
    \label{fig:adapt_1366}
\end{figure}

\begin{figure}[H]
    \centering
    \includegraphics[width=1\textwidth]{figs/360+800.png}
    \caption{Тестирование адаптивности для разрешения $360\times800$}
    \label{fig:adapt_360}
\end{figure}

\section{Тестирование верстки}

Проверка валидности верстки выполнена с использованием официального сервиса W3C Markup Validation Service. В ходе анализа исходного кода веб-страницы были выявлены предупреждения и информационные сообщения, касающиеся структуры документа. Фрагмент отчета о валидации представлен на \cref{fig:validator}, поскольку весь отчет занимает непозволительно большой объем.

\begin{figure}[H]
    \centering
    \includegraphics[height=\textwidth, page=1, angle=-90]{figs/validator.pdf}
    \caption{Результаты проверки верстки в W3C Validator}
    \label{fig:validator}
\end{figure}



\section{Технический аудит}

Технический аудит сайта проводился с использованием инструмента W3Techs. В результате анализа определен стек технологий, применяемых на сервере и клиенте. Результаты аудита представлены на \cref{fig:w3techs_audit}.

\begin{figure}[H]
    \centering
    \includegraphics[height=\textwidth, page=1]{figs/w3techs.pdf}
    \caption{Результаты технического аудита веб-приложения (W3Techs)}
    \label{fig:w3techs_audit}
\end{figure}

\section{Тестирование сетевых подключений}

Тестирование сетевых подключений проводилось с использованием встроенных инструментов разработчика (DevTools, вкладка Network). Общий каскад загрузки компонентов веб-страницы представлен на \cref{fig:network_cascade}.

\begin{figure}[H]
    \centering
    \includegraphics[width=\textwidth]{figs/network_cascade.png}
    \caption{Каскад загрузки веб-страницы}
    \label{fig:network_cascade}
\end{figure}

Для детального анализа выбран сетевой запрос к основному документу (\textit{www.digital-energy.ru}). Во вкладке Headers (\cref{fig:network_header}) определены базовые параметры подключения:
\begin{itemize}
    \item Тип подключения: HTTPS.
    \item Метод запроса: GET.
    \item Код состояния: 200 (успешное выполнение).
\end{itemize}

\begin{figure}[H]
    \centering
    \includegraphics[width=\textwidth]{figs/network_header.png}
    \caption{Заголовки запроса и ответа (Headers)}
    \label{fig:network_header}
\end{figure}

Во вкладке Response (\cref{fig:network_responce}) зафиксировано тело ответа, представляющее собой исходный HTML-код загружаемого документа.

\begin{figure}[H]
    \centering
    \includegraphics[width=\textwidth]{figs/network_response.png}
    \caption{Тело ответа сервера (Response)}
    \label{fig:network_responce}
\end{figure}

Детальная информация о параметрах маршрутизации и дополнительных заголовках показана на \cref{fig:network_www}.

\begin{figure}[H]
    \centering
    \includegraphics[width=\textwidth]{figs/network_www.png}
    \caption{Дополнительная информация о сетевом запросе}
    \label{fig:network_www}
\end{figure}

Анализ файлов cookie, передаваемых в рамках сессии, представлен на \cref{fig:network_cookie}. Наличие флага Secure указывает на то, что файлы cookie используют защищенное соединение.

\begin{figure}[H]
    \centering
    \includegraphics[width=\textwidth]{figs/network_cookie.png}
    \caption{Параметры файлов cookie}
    \label{fig:network_cookie}
\end{figure}

\section{Определение элементов автоматизации тестирования на основе локаторов}

В рамках подготовки к автоматизированному тестированию был исследован исходный код веб-элемента (логотипа компании) методом «белого ящика» с использованием вкладки Elements (\cref{fig:logo_code}).

\begin{figure}[H]
    \centering
    \includegraphics[width=\textwidth]{figs/logo_code.png}
    \caption{Исходный код элемента во вкладке Elements}
    \label{fig:logo_code}
\end{figure}

Фрагмент программного кода исследуемого элемента:
\begin{listing}[H]
\caption{Исходный HTML-код веб-элемента логотипа}
\label{lst:logo_code}
\begin{minted}[breaklines]{html}
<a id="logo" href="https://www.digital-energy.ru" data-supplied-ml-starting-dark="true" data-supplied-ml-starting="false" data-supplied-ml="false">
    <img class="stnd skip-lazy default-logo" width="250" height="75" alt="Ассоциация &quot;Цифровая Энергетика&quot;" src="https://www.digital-energy.ru/wp-content/uploads/2021/02/de-logo-250-color.png" srcset="https://www.digital-energy.ru/wp-content/uploads/2021/02/de-logo-250-color.png 1x, https://www.digital-energy.ru/wp-content/uploads/2021/02/de-logo-500-color.png 2x">
</a>
\end{minted}
\end{listing}

Элемент сформирован на основе тега \textit{a}, внутрь которого вложен тег \textit{img}. Ключевые атрибуты тега \textit{a}, которые можно использовать для поиска элемента в DOM-дереве: \textit{id="logo"} и \textit{href="https://www.digital-energy.ru"}. 

Варианты селекторов для обращения к данному элементу:
\begin{itemize}
    \item CSS-локатор по идентификатору: \textit{\#logo}.
    \item XPath-локатор: \textit{//a[@id='logo']}.
\end{itemize}


Для расширения тестового покрытия были дополнительно проанализированы 4 элемента графического интерфейса. Результаты исследования исходного кода элементов представлены на \cref{fig:calendar_button,fig:hamburger,fig:search,fig:archive}.

\begin{figure}[H]
    \centering
    \includegraphics[width=0.9\textwidth]{figs/calendar_button.png}
    \caption{Исходный код элемента «Календарь мероприятий»}
    \label{fig:calendar_button}
\end{figure}

Кнопка «Календарь мероприятий» реализована через тег \textit{a} с набором стилевых классов.
\begin{itemize}
    \item CSS-локатор: \textit{a.nectar-button.large.regular}
    \item XPath-локатор: \textit{//a[contains(text(), 'КАЛЕНДАРЬ МЕРОПРИЯТИЙ')]}
\end{itemize}

\begin{figure}[H]
    \centering
    \includegraphics[width=0.9\textwidth]{figs/hamburger.png}
    \caption{Исходный код элемента мобильного меню}
    \label{fig:hamburger}
\end{figure}

Элемент управления боковым меню определен тегом \textit{a} с классом \textit{slide-out-widget-area-toggle}.
\begin{itemize}
    \item CSS-локатор: \textit{a.slide-out-widget-area-toggle}
    \item XPath-локатор: \textit{//a[contains(@class, 'slide-out-widget-area-toggle')]}
\end{itemize}

\begin{figure}[H]
    \centering
    \includegraphics[width=0.9\textwidth]{figs/search.png}
    \caption{Исходный код элемента «Поиск»}
    \label{fig:search}
\end{figure}

Иконка вызова формы поиска вложена в элемент списка \textit{li} с уникальным идентификатором.
\begin{itemize}
    \item CSS-локатор: \textit{\#search-btn a}
    \item XPath-локатор: \textit{//li[@id='search-btn']/a}
\end{itemize}

\begin{figure}[H]
    \centering
    \includegraphics[width=0.9\textwidth]{figs/archive.png}
    \caption{Исходный код ссылки «Архив новостей»}
    \label{fig:archive}
\end{figure}

Ссылка на архив новостей реализована стандартным тегом \textit{a} с указанием относительного пути.
\begin{itemize}
    \item CSS-локатор: \textit{a[href*='digital-energy.ru/news']}
    \item XPath-локатор: \textit{//a[contains(text(), 'Архив новостей')]}
\end{itemize}


\section{Автоматизация тестирования}

Для автоматизации процесса тестирования применялось браузерное расширение Selenium IDE. Был разработан сценарий дымового тестирования, включающий последовательный переход по ключевым разделам веб-приложения.Сценарий автоматизированного тестирования зафиксирован в табличной форме (\cref{tab:selenium_steps}).

\begin{longtblr}[
    caption = {Сценарий автоматизированного тестирования},
    label = {tab:selenium_steps},
]{
    colspec = { Q[l,m] X[l,m] X[l,m] },
    hlines, vlines,
    rowhead = 1,
}
    № п.п. & Последовательность действий & Ожидаемый результат \\
    1 & Открыть веб-приложение по адресу \url{https://www.digital-energy.ru/} & Приложение открыто на главной странице \\
    2 & Установить размер окна браузера $1936\times1096$ & Размер окна успешно изменен \\
    3 & Нажать на пункт меню «Об ассоциации» & Осуществлен переход на соответствующую страницу \\
    4 & Навести мышь на «Деятельность» & Открылось всплывающее меню \\
    5 & Нажать на пункт меню «Новости» & Осуществлен переход в раздел <<Образование>> \\
    6 & Нажать на пункт меню «Новости» & Осуществлен переход в раздел новостей \\
    7 & Нажать на кнопку «Больше новостей» & Загружен дополнительный блок контента \\
    8 & Нажать на иконку вызова бокового меню & Открыто боковое меню навигации \\
    9 & Нажать на пункт «Учредители» в боковом меню & Осуществлен переход на страницу учредителей \\
    10 & Нажать на логотип в шапке сайта & Осуществлен возврат на главную страницу \\
    11 & Нажать на пункт меню «Контакты» & Осуществлен переход в раздел контанктов \\
\end{longtblr}

Визуальное представление записанных команд в интерфейсе Selenium IDE приведено на \cref{fig:selenium_script}.

\begin{figure}[H]
    \centering
    \includegraphics[width=\textwidth]{figs/selenium_script.png}
    \caption{Сценарий автоматизированного тестирования в Selenium IDE}
    \label{fig:selenium_script}
\end{figure}

При воспроизведении сценария все шаги были выполнены успешно. Лог выполнения с подтверждением корректности прохождения каждого этапа представлен на \cref{fig:selenium_log}.

\begin{figure}[H]
    \centering
    \includegraphics[width=\textwidth]{figs/selenium_log.png}
    \caption{Лог успешного выполнения автоматизированного теста}
    \label{fig:selenium_log}
\end{figure}

Экспортированный исходный код сценария в формате JSON приведен в \cref{lst:selenium_code}.

\section{Применение искусственного интеллекта в тестировании}

В рамках выполнения проекта для оптимизации процессов проектирования тестовых мероприятий использовалась большая языковая модель. Были сформулированы специализированные промпты для решения трех различных задач.

\subsection{Запрос на предоставление определения из сферы тестирования}
Нейросети была поставлена задача сформулировать академическое определение нагрузочного тестирования в контексте веб-приложений. Результат выполнения запроса представлен на \cref{fig:ai_def}.

\begin{figure}[H]
    \centering
    \includegraphics[width=\textwidth]{figs/ai_def.png}
    \caption{Генерация определения нагрузочного тестирования}
    \label{fig:ai_def}
\end{figure}

\subsection{Запрос на составление вопроса по техникам тестирования}
Для проверки образовательных возможностей ИИ был сформирован запрос на создание тестового вопроса по техникам тест-дизайна (эквивалентное разбиение, анализ граничных значений). Сгенерированный вопрос с вариантами ответа показан на \cref{fig:ai_question}.

\begin{figure}[H]
    \centering
    \includegraphics[width=\textwidth]{figs/ai_question.png}
    \caption{Генерация тестового вопроса по тест-дизайну}
    \label{fig:ai_question}
\end{figure}

\subsection{Запрос на составление тест-кейсов}
С целью автоматизации рутинной работы тестировщика, нейросети был передан URL исследуемого веб-приложения с требованием сгенерировать 5 базовых тест-кейсов. Результат, оформленный в табличном виде, зафиксирован на \cref{fig:ai_cases}.

\begin{figure}[H]
    \centering
    \includegraphics[width=\textwidth]{figs/ai_cases.png}
    \caption{Генерация базовых тест-кейсов для тестируемого ресурса}
    \label{fig:ai_cases}
\end{figure}
